\begin{remark}
\label{remark-category-convention-boundary}
\begin{reference}
\cite[Chapter I, Section 1.1, p. 122]{Tohoku}
\end{reference}
The cited source builds three conventions into the word category: the class
of objects is nonempty, it may be a proper class, and the sets of morphisms
for distinct ordered pairs of objects are disjoint. Definition
\ref{definition-category} instead permits the empty category and uses a set
of objects by default; the selected exceptions to the latter convention are
listed in Remark \ref{remark-big-categories}. The discussion following
Definition \ref{definition-category} treats disjointness of the morphism
sets as optional. Apart from these three convention choices, the data and
axioms are the same.

The disjointness convention changes no categorical structure. Given a
category $\mathcal{C}$, let $\mathcal{C}^\dagger$ have the same objects and
put
$$
\Mor_{\mathcal{C}^\dagger}(x, y) =
\{(x, f, y) \mid f \in \Mor_\mathcal{C}(x, y)\}.
$$
Define composition and identities by
$$
(y, g, z) \circ (x, f, y) = (x, g \circ f, z), \qquad
\text{id}^{\dagger}_x = (x, \text{id}_x, x).
$$
Associativity and the identity axioms follow from those of $\mathcal{C}$.
The displayed maps define a functor
$T : \mathcal{C} \to \mathcal{C}^\dagger$ which is the identity on objects
and sends $f \in \Mor_\mathcal{C}(x, y)$ to $(x, f, y)$. The identity on
objects together with $(x, f, y) \mapsto f$ defines an inverse functor.
Thus pairwise disjoint morphism sets can always be imposed by tagging arrows
with their source and target; object size and nonemptiness remain convention
choices.
\end{remark}
